\chapter{INTRODUCTION}
\label{ch:introduction}

\todoresolve{REVIEW NARRATIVE ARC AND CHAPTER TRANSITIONS, ESPECIALLY PLACEMENT OF CHAPTER 6 (BOOSTING INFERENCE).}

\begin{aigenblock}

\section{Statistical and Computational Costs in Sequential Learning}

Many problems in science and engineering require learning from observations that arrive sequentially over time. A robotic controller selects motor commands and observes the resulting motion of a physical system. A recommendation platform observes how users respond to suggested items and adapts future recommendations as more responses arrive. An epidemiological model assimilates weekly reported case counts to infer the unobserved spread of an infectious disease and evaluate candidate interventions. Across these settings, a learning system faces two distinct resource constraints: the statistical cost of collecting informative data from the environment, and the computational cost of representing, updating, and optimizing models, beliefs, and policies once those data are in hand.

In sequential decision making, actions play a dual role because they determine both immediate outcomes and the observations available for future learning. A policy trained only on historical logs may have little evidence about actions that the logging policy rarely selected. Collecting new trajectories through online interaction can supply the missing information, yet every physical trial takes time and incurs real-world cost. Even after data are collected, turning past observations into a decision or a parameter estimate requires computation. Evaluating a policy class whose description grows with past value estimates, optimizing over continuous multimodal action distributions, quantifying uncertainty across sequentially constructed tree ensembles, integrating over unobserved state trajectories with particle filters, or branching over future observation histories in a belief tree can each become the primary bottleneck in practice.

Analyzing statistical sample complexity while treating computation and representation as exact, or designing fast heuristics without tracking the approximation errors they introduce, can obscure the dominant source of error. A statistically optimal reinforcement learning method that stores an ever-growing sum of past value functions or solves an intractable nonconvex program at every episode cannot be deployed at scale. Conversely, replacing an exact policy update, an exact posterior belief, or an exact lookahead tree with a tractable approximation introduces projection bias, Monte Carlo variance, or structural lookahead error that directly enters the performance guarantee. The central question of this dissertation is how the representation and use of information determine both the statistical cost of data collection and the computational cost of inference and decision making.

\section{Overview of the Dissertation}

This dissertation is organized around seven research chapters that progress from learning policies under observed states and contexts, to statistical inference for sequentially constructed regression ensembles, and finally to parameter estimation and control when the underlying state is only partially observed. Throughout the dissertation, we number algorithms by chapter and state the exact assumptions, sample-size regimes, and approximation errors attached to each theoretical guarantee.

\subsection{Learning from Offline and Online Data}

We begin in Chapter~\ref{ch:hybrid-rl} with the most direct question about combining information sources in reinforcement learning: what can an existing offline dataset tell us, and what must an agent still learn through online interaction? This chapter is based on joint work with Ziping Xu \citep{tan2024natural} and joint work with Wei Fan and Yuting Wei \citep{tan2024hybrid}.

Purely offline reinforcement learning avoids new environment interaction by fitting a pessimistic value function to logged transitions, but standard guarantees require the behavior policy to cover the state--action distribution of an optimal policy. When the offline dataset lacks full coverage on part of the state--action space, purely offline methods can fail arbitrarily badly, whereas purely online methods ignore the logged transitions altogether and explore the entire state--action space from scratch. Existing hybrid reinforcement learning algorithms either assume that the offline dataset already satisfies single-policy concentrability or require knowing an explicit partition of the state--action space in advance.

In the first part of Chapter~\ref{ch:hybrid-rl}, we study a natural extension of optimistic online reinforcement learning under general function approximation (Algorithm~\ref{alg:ch02-alg1}), which simply appends the offline dataset to the empirical Bellman-error constraints used to construct the confidence set. Without knowing where the offline data are informative, the analysis partitions the state--action space into a region where performance is controlled by a partial offline concentrability coefficient and a complementary region where performance is controlled by an online Bellman eluder dimension. Taking the infimum over all such analytical partitions shows that optimism automatically directs online exploration toward regions where offline coverage is weak.

In the second part of Chapter~\ref{ch:hybrid-rl}, we ask whether combining offline and online data can strictly improve upon the minimax sample complexity barriers of both purely offline and purely online learning in linear Markov decision processes (MDPs), where transition kernels and rewards are linear in a known $d$-dimensional feature map $\phi_h(s,a) \in \R^d$ over horizon $H$. We develop two computationally tractable algorithms targeting different evaluation criteria. Algorithm~\ref{alg:ch02-alg2} takes an online-to-offline perspective for final-policy probably approximately correct (PAC) sample complexity: it uses reward-agnostic online exploration to collect samples in feature directions that are poorly covered by the offline covariance matrix, and then fits a variance-weighted pessimistic value function on the combined dataset. This replaces the offline information matrix $\Sigma_{\mathrm{off},h}^*$ in the instance-dependent minimax suboptimality bound with the combined matrix $\Sigma_{\mathrm{off},h}^* + \Sigma_{\mathrm{on},h}^*$ without requiring full-rank offline coverage. Algorithm~\ref{alg:ch02-alg3} takes an offline-to-online perspective for cumulative regret over $T$ online episodes given $N_{\mathrm{off}}$ offline episodes, achieving a high-probability regret bound of order
\begin{equation}
\label{eq:intro-hybrid-regret}
\Reg(T) \le \widetilde{O}\!\left( \inf_{\mathcal{X}_{\mathrm{off}},\,\mathcal{X}_{\mathrm{on}}} \left\{ c_{\mathrm{off}} T \sqrt{\frac{d H^3}{N_{\mathrm{off}}}} + \sqrt{d_{\mathrm{on}} d H^3 T} \right\} \right)
\end{equation}
once $T$ and $N_{\mathrm{off}}$ exceed the large-sample threshold $\widetilde{\Omega}(d^{13} H^{14})$, where $c_{\mathrm{off}}$ measures partial offline coverage on $\mathcal{X}_{\mathrm{off}}$ and $d_{\mathrm{on}} \le d$ is the dimension of the subspace spanned by $\mathcal{X}_{\mathrm{on}}$. When the offline dataset is large and covers all but a $d_{\mathrm{on}}$-dimensional subspace of features, this bound improves upon the online minimax rate $\widetilde{\Theta}(\sqrt{d^2 H^3 T})$ by replacing one factor of $\sqrt{d}$ with $\sqrt{d_{\mathrm{on}}}$.

\subsection{Learning Shared Structure from Heterogeneous Data}

Chapter~\ref{ch:latent-bandits} examines a second way historical data can reduce the cost of online learning, based on joint work with Chinmaya Kausik and Ambuj Tewari \citep{kausik2024leveraging}. Rather than supplying additional transitions from the same MDP, historical data in recommendation and personalization systems often consist of logged interactions from many heterogeneous users or tasks, each governed by its own unknown reward parameter $\theta_i \in \R^d$. When the user parameters lie near a shared $r$-dimensional subspace with $r \ll d$, offline data across users can identify this lower-dimensional representation so that online learning for a new user only needs to explore an $r$-dimensional space.

The central obstacle is that historical logs for each user are collected by unknown, potentially history-dependent behavior policies over very short trajectories, sometimes as short as two actions per user. Under adaptive action selection and non-spherical feature covariances, naive empirical second-moment matrices of feature--reward products are biased by within-trajectory dependencies and observation noise, and existing meta-learning or representation-learning guarantees require trajectory lengths to grow polynomially with the ambient dimension $d$.

We construct a moment-corrected spectral estimator (Algorithm~\ref{alg:sold}) that pairs two observations per offline user through importance-weighted cross-moments to cancel the confounding from adaptive action selection and observation noise. We prove that the subspace estimation error in the sine-theta distance decays at the rate $\widetilde{O}(\sqrt{d / N_{\mathrm{off}}})$ in the number of offline users $N_{\mathrm{off}}$ even with a constant number of observations per user. We then pass the estimated $r$-dimensional subspace to two online contextual bandit algorithms for a new user: a subspace-constrained confidence-set bandit method (Algorithm~\ref{alg:local-ucb}) and a computationally efficient projected confidence-ball variant (Algorithm~\ref{alg:proball-ucb}), both achieving cumulative regret of order $\widetilde{O}(r\sqrt{T} + T\sqrt{d/N_{\mathrm{off}}})$.

\subsection{Statistically and Computationally Efficient Actor--Critic Methods}

Chapters~\ref{ch:hybrid-rl} and~\ref{ch:latent-bandits} rely on value-based or confidence-set policies that select actions by maximizing an optimistic value estimate. In Chapter~\ref{ch:actor-critic}, we turn to actor--critic methods, which maintain an explicit policy (the actor) alongside a value estimate (the critic). This chapter is based on joint work with Wei Fan and Yuting Wei \citep{tan2025actor} and joint work with Max Qiushi Lin, Reza Asad, Haque Ishfaq, Csaba Szepesv\'ari, and Sharan Vaswani \citep{lin2026parametric}.

Actor--critic methods are widely used in practice because incremental policy updates are convenient to implement, yet existing theoretical guarantees under general function approximation either require global coverage assumptions or suffer suboptimal sample complexity exponents such as $\widetilde{O}(T^{2/3})$ or $\widetilde{O}(T^{3/4})$ regret. In the first part of Chapter~\ref{ch:actor-critic}, we identify the exact statistical mechanism behind this gap. When the actor performs natural multiplicative weight updates with step size $\eta$ using a sequence of critics $\widehat{Q}_1, \dots, \widehat{Q}_t$, unrolling the recursion from a uniform initial policy yields
\begin{equation}
\label{eq:intro-ac-unroll}
\pi_{t+1}(a \mid s) \propto \exp\!\left( \eta \sum_{\tau=1}^t \widehat{Q}_\tau(s,a) \right).
\end{equation}
Each new critic update adds another function from the critic class $\mathcal{F}$ to the exponent, so after $K$ distinct critic updates the log-covering number of the induced policy class $\Pi_T$ grows linearly with $K$. If the critic is updated at every episode ($K = T$), uniform concentration over the policy-dependent Bellman target class breaks down and yields a vacuous linear regret bound unless $\Pi_T$ happens to have a small cover by special structure (the case treated by Algorithm~\ref{alg:DOUHUA}).

One might hope to control $K$ simply by updating the critic rarely, only when its empirical Bellman loss exceeds a confidence threshold. However, if a frozen critic targets the current policy value $Q^{\pi_t}$, the regression target $\mathcal{T}_h^{\pi_t} \widehat{Q}_{h+1}$ drifts whenever the actor updates $\pi_t$, which destroys optimism and forces frequent critic switches anyway. Algorithm~\ref{alg:NORA} resolves this obstacle with two coupled modifications: the critic targets the optimal action-value function $Q^*$ via the Bellman optimality operator $(\mathcal{T}_h f)(s,a) = \E[r_h + \max_{a'} f(s',a') \mid s,a]$ so the target operator no longer depends on the changing actor, and whenever the rare-switching test triggers a critic refit, the actor resets to a uniform distribution before resuming mirror-descent updates toward the new critic. Because the maximum operator is $1$-Lipschitz in the supremum norm, the target class covering number is controlled by $\mathcal{F}$ alone, while uniform resets cap the Kullback--Leibler divergence to the new greedy policy at $\log|\mathcal{A}|$ on each of the $K = O(d_{\mathrm{BE}} H \log T)$ switching intervals. This achieves $\widetilde{O}(\sqrt{T})$ regret and $\widetilde{O}(\epsilon^{-2})$ sample complexity under general function approximation without global coverage.

Although Algorithm~\ref{alg:NORA} is statistically optimal in $\epsilon$, evaluating an implicit policy of the form~\eqref{eq:intro-ac-unroll} still requires summing over all critics accumulated within an interval, and global optimistic confidence-set optimization remains computationally intractable. The second part of Chapter~\ref{ch:actor-critic} addresses both computational bottlenecks in linear MDPs (Algorithms~\ref{alg:optimistic_actor_critic}--\ref{alg:lmc}). Instead of storing past critics, Algorithm~\ref{alg:actor_npg} parameterizes the actor with fixed-dimensional log-linear logits $z_w(s,a) = \langle w, \psi(s,a) \rangle$ and approximates the ideal logit update $z_{w_t} + \eta \widehat{Q}_t$ via weighted least-squares regression on a compact $G$-optimal design coreset of policy features, while Algorithm~\ref{alg:lmc} replaces confidence-set optimization in the critic with noisy gradient updates via Langevin Monte Carlo. This keeps per-step policy evaluation cost constant in $T$ and yields $\widetilde{O}(\epsilon^{-4})$ sample complexity under on-policy sampling and $\widetilde{O}(\epsilon^{-2})$ sample complexity when off-policy critic data reuse is permitted.

\subsection{Learning Continuous-Action Diffusion Policies}

In Chapter~\ref{ch:diffusion-policies}, based on joint work with Zhihan Huang and Yuting Wei \citep{tan2026learning}, we extend the study of policy representations to continuous action spaces $\mathcal{A} \subseteq \R^{d_a}$, where optimal or near-optimal policies may be multimodal and poorly approximated by unimodal Gaussian or log-linear families. Diffusion models represent complex continuous distributions by transforming an initial Gaussian noise draw $z_K \sim \mathcal{N}(0, I_{d_a})$ through $K$ reverse-time denoising transitions $z_K \to z_{K-1} \to \dots \to z_0 = a$ conditioned on the state $s$.

While expressive, diffusion policies introduce a fundamental obstacle for reinforcement learning: the terminal action density $\pi_\theta(a \mid s) = \int p_\theta(z_{0:K} \mid s)\,dz_{1:K}$ integrates over all unobserved denoising paths and cannot be evaluated in closed form, so standard policy-gradient and natural-policy-gradient updates cannot be applied directly to $\log \pi_\theta(a \mid s)$. Chapter~\ref{ch:diffusion-policies} develops and analyzes three algorithmic routes that bypass this intractability in different ways:
\begin{enumerate}
\item \textbf{Vanilla diffusion policy gradient (Algorithm~\ref{alg:dppo-vanilla-pg}):} Using either an unbiased denoising-path score estimate for stochastic denoising diffusion probabilistic model (DDPM) samplers or the exact endpoint score for deterministic denoising diffusion implicit model (DDIM) samplers, we establish $O(T^{3/4})$ expected regret under reset or fixed-state mixing access, up to comparison-policy and compatible-approximation errors.
\item \textbf{Monte Carlo diffusion natural policy gradient (Algorithm~\ref{alg:dppo-reinforce}):} Under a deterministic DDIM flow, we evaluate the exact endpoint score via the change-of-variables formula and perform natural policy gradient updates under reset, mixing, or hybrid logged-trajectory access, achieving $O(T^{3/4})$ high-probability regret up to comparison-policy and advantage-transfer errors without explicit density-ratio dependence.
\item \textbf{Diffusion policy behavior cloning of an optimistic critic (Algorithm~\ref{alg:dipo-bc}):} To avoid reset-coverage assumptions in finite-horizon reinforcement learning, we construct an optimistic action-value critic (or backward fitted-$Q$ critic with offline data) and train the diffusion policy by denoising score matching on the critic's greedy actions, achieving $\widetilde{O}(\sqrt{T})$ regret up to cumulative score-matching and sampler errors.
\end{enumerate}

\subsection{Statistical Inference for Boosted Regression and Additive Models}

Every method in Chapters~\ref{ch:hybrid-rl} through~\ref{ch:diffusion-policies} relies on sequential function approximation to estimate values or policies, where each iterate depends on the residuals or targets produced by earlier iterates. In Chapter~\ref{ch:boosting-inference}, based on joint work with Haimo Fang and Giles Hooker \citep{fang2025statistical} and joint work with Haimo Fang, Jonathan Pipping, and Giles Hooker \citep{fang2026statistical}, we examine this sequential function-fitting dynamics directly in supervised regression, asking how to conduct rigorous frequentist statistical inference---confidence intervals, prediction intervals, and hypothesis tests---for gradient boosted decision trees and explainable boosting machines.

Standard gradient boosting fits each new regression tree to the empirical residuals of all preceding trees. Because every tree depends on the full sequence of earlier adaptive splits and leaf values, the final ensemble is a highly coupled functional of the response vector, and classical $U$-statistic arguments used to establish asymptotic normality for random forests do not apply. The moving-average regularization of \citet{zhou2022boulevard} achieves a central limit theorem (CLT) by rescaling the ensemble at step $b$ by $(1-\lambda)/b + \lambda$, which causes the ensemble to converge to the fixed point of a contraction mapping. However, that moving average shrinks the asymptotic prediction to $\mS(\mI + \frac{1-\lambda}{\lambda}\mS)^{-1}\vy$, equivalent to kernel ridge regression with penalty $(1-\lambda)/\lambda > 0$, which introduces substantial shrinkage bias toward zero and degrades predictive accuracy unless $\lambda=1$.

In the first part of Chapter~\ref{ch:boosting-inference}, we show that integrating moving-average regularization with either random tree dropout (Algorithm~\ref{alg:random-dropout}) or parallel tree growing with structured dropout (Algorithm~\ref{alg:structured-dropout}) changes the contraction fixed point itself. Increasing the dropout probability $p$ or the number of parallel trees $K$ weakens the effective ridge penalty on the signal, allowing the ensemble to interpolate smoothly between regularized boosting and random forests while preserving an explicit asymptotic normal distribution for point predictions, prediction intervals, and chi-squared variable-importance tests. In the second part of Chapter~\ref{ch:boosting-inference}, we extend this moving-average contraction framework to cyclic univariate boosting in Generalized Additive Models (Algorithm~\ref{alg:brebm_ident}), establishing feature-wise asymptotic normality that equips glass-box additive shape functions $f_k(x_k)$ with pointwise confidence bands and formal hypothesis tests for feature significance and nonlinearity.

\subsection{Transition to Partial Observability: Parameter Estimation and Control}

In Chapters~\ref{ch:hybrid-rl} through~\ref{ch:diffusion-policies}, the decision maker observes the Markov state $s_h$ or context $x_t$ directly at each step, so the representation problem concerns how to parameterize value functions and policies over an observed state space. In many scientific and physical domains, however, the underlying Markov state $X_t$ is hidden and can only be inferred through a sequence of noisy, incomplete observations $Y_{1:t}$. As a simple physical illustration, consider applying a braking force to stop a moving cart when the sensor measures only noisy position and not velocity: two carts at the exact same measured position moving in opposite directions require opposite braking forces, so no memoryless policy of the current observation can control the system. Inferring velocity requires tracking changes in position across time and maintaining a posterior belief distribution over the latent state.

Chapters~\ref{ch:pomp-estimation} and~\ref{ch:pomdp-control} make hidden-state inference itself part of the statistical and computational problem. Sequential Monte Carlo (particle filtering) approximates the latent state distribution using an ensemble of $J$ weighted particles simulated forward from the transition dynamics and resampled according to observation likelihoods. While particle filters are plug-and-play---requiring only a forward simulator for the latent dynamics rather than closed-form transition densities---their discrete resampling step and finite particle budget $J$ create distinct obstacles for parameter estimation and for online control. Chapter~\ref{ch:pomp-estimation} studies the statistical estimation problem for passive observation sequences, and Chapter~\ref{ch:pomdp-control} brings particle belief tracking into online action selection and planning.

\subsubsection{Parameter Estimation under Partial Observability}

Chapter~\ref{ch:pomp-estimation}, based on joint work with A.~J.~Abkemeier, Kunyang He, Giles Hooker, and Edward L.~Ionides \citep{tan2024accelerated}, studies maximum likelihood and Bayesian parameter estimation in partially observed Markov processes (POMPs). Because a bootstrap particle filter produces an unbiased likelihood estimate using only a forward simulator and observation density evaluations, differentiating the filter via automatic differentiation (AD) is an attractive route to gradient-based parameter estimation. The obstacle is that discrete ancestor resampling changes discontinuously when the parameter $\theta$ is perturbed under fixed random numbers, making the raw particle-filter log-likelihood estimate a step function with zero or undefined pathwise gradient.

Existing attempts to differentiate through particle filters either stop gradients at the resampling weights \citep{naesseth18,jonschkowski18}, which discards past weight sensitivities and incurs $O(1)$ asymptotic score bias, or retain full unrescaled ancestral weight ratios \citep{poyiadjis11,scibior21}, whose score variance grows linearly or quadratically with the time horizon $N$ and typically requires evaluating latent transition densities. We introduce a smooth, plug-and-play off-parameter particle filter family parameterized by a weight discount factor $\alpha \in [0,1]$ (Algorithm~\ref{alg:mop}) and its single-pass differentiable counterpart at the baseline parameter (Algorithm~\ref{alg:dmop}). By coupling nearby parameter evaluations $\theta$ to a baseline filter run at $\phi$ with fixed resampling ancestors and discounting carried weight ratios by $w_{n,j}^{P,\theta} = (w_{n-1,j}^{F,\theta})^\alpha$ before applying an importance-weight correction at step $n$, Algorithm~\ref{alg:dmop} yields a pathwise score estimator that recovers the biased stop-gradient filter at $\alpha = 0$ and the consistent high-variance filter as $\alpha = 1$ and $J \to \infty$, while achieving an explicit bias--variance tradeoff for $\alpha \in (0,1)$ with squared bias $O(N\alpha^{2/(1-\alpha)})$ and particle variance $O(N(1+\log J)^2/(J(1-\alpha)^2))$ under uniform mixing assumptions.

Because stochastic gradient ascent with particle-filter gradients converges rapidly near a local mode but can stall in rough likelihood regions far from the maximum likelihood estimate, Algorithm~\ref{alg:ifad} combines an iterated filtering warm start \citep{ionides15} with subsequent discounted AD gradient optimization. We prove finite-sample and asymptotic convergence guarantees for this two-stage procedure and demonstrate substantial empirical improvements on the nonlinear Dhaka cholera model of \citet{king08} and the 38-parameter \textit{Daphnia} ecological panel model of \citet{yang25daphnia}.

\subsubsection{Control under Partial Observability}

Finally, Chapter~\ref{ch:pomdp-control}, based on joint work with Ziping Xu \citep{tan2026pomdp}, studies online reinforcement learning and planning in partially observable Markov decision processes (POMDPs). Even when a latent dynamics model $\bar{\theta}_k$ is estimated online from interaction trajectories, computing an optimal action requires planning in belief space over future action and observation branches. Exact belief-tree search branches over all possible actions and discretized observations at every lookahead step, incurring computational complexity exponential in the planning horizon, whereas the classical QMDP approximation $Q_{\mathrm{QMDP}}(b,a) = \E_{s \sim b}[Q_{\mathrm{MDP}}^*(s,a)]$ averages the fully observed state--action value over current belief particles in $O(J)$ time but ignores the value of future information gathering because it assumes the latent state will become fully observed after the next transition.

Algorithm~\ref{alg:pf_mcts} bridges this gap by combining epoch-wise maximum likelihood model estimation, bootstrap particle filtering with $J$ particles, and a depth-$d$ Monte Carlo belief-tree planner (Algorithm~\ref{alg:mctq}) with $N$ rollouts per node that expands observation branches for $d$ steps before evaluating leaf beliefs with the fully observed state--action value function $Q_{\bar{\theta}_k}^{*,S}$. Under uniform Doeblin mixing and bounded positive observation densities, we prove an end-to-end cumulative regret bound over $T$ steps of the form
\begin{equation}
\label{eq:intro-pomdp-regret}
\frac{\Reg(T)}{T} \lesssim \widetilde{O}(T^{-1/4}) + C_{\mathrm{pf}} J^{-1/(2\kappa_{\mathrm{pf}})} + C_{\mathrm{mc}}(d,M) N^{-1/2} + C \gamma^{d+1} \overline{\Xi}_T,
\end{equation}
which explicitly separates the statistical interaction error $\widetilde{O}(T^{-1/4})$ of online model learning from the three computational approximations: particle belief tracking error $O(J^{-1/(2\kappa_{\mathrm{pf}})})$, Monte Carlo tree rollout error $O(C_{\mathrm{mc}}(d,M) N^{-1/2})$ across $M$ observation bins, and structural lookahead error $O(\gamma^{d+1}\overline{\Xi}_T)$ from deferring the full-observability leaf bootstrap by $d$ steps. We validate each of the four scaling rates separately on a mixing POMDP that satisfies the theorem's assumptions exactly, and we instantiate the depth-zero particle state-rollout planner with learned dynamics and actor--critic networks on continuous-control and mobile-health benchmarks.

Chapter~\ref{ch:conclusion} concludes the dissertation by synthesizing the common principles connecting statistical coverage, structural representation, and computational approximation across Chapters~\ref{ch:hybrid-rl} through~\ref{ch:pomdp-control} and outlining open directions for future research.
\end{aigenblock}
